% !TEX root = ../main.tex
%% Threats to Validity

\section{Threats to Validity}%
\label{sec:threats}

\subsection{Construct Validity}

The primary construct validity concern is the ground-truth reference
model used for F1 scoring and JSD/Frobenius measurement.  On the AV
CPS benchmark, the ground-truth state set and transition set were
derived directly from the natural-language requirements specification,
with states explicitly enumerated in the text.  This limits annotation
subjectivity but also means that the benchmark is structurally
self-contained: all ground-truth information is available to the oracle
during model construction.  The JSD and Frobenius metrics compare
NeSy-MBST output against a reference matrix whose transition
probabilities were randomly initialised from the ground-truth topology
(seed~42), not from empirically measured operational profiles.  These
metrics therefore quantify calibration improvement over a random
baseline rather than fidelity to real-world usage distributions.

The single-author setting introduces annotation bias risk for
ground-truth labelling.  On the AV benchmark this risk is mitigated by
the formal specification source; on the e-commerce models the
ground-truth is identical to the framework's own input, making F1
trivially~1.0 by constructi is a limitation acknowledged in
Section~\ref{sec:eval:ablation}.

\subsection{Internal Validity}

The L\* oracle's consistency is not formally guaranteed for
three-valued responses.  In our experiments the \texttt{Unsure}
escalation path was never triggered on the AV benchmark (oracle
resolution rate: 94\% direct, 6\% escalated to SUT execution), so
consistency held empirically.  Benchmarks with more ambiguous
requirements may produce higher escalation rates; systematic
inconsistency would degrade convergence.

The ablation uses a fixed random seed (42) for the reference matrix and
the test generator.  Results for a single seed are reproducible but may
not represent the distribution of outcomes across seeds.  The
implementation is publicly available; practitioners can re-run
\texttt{run\_ablation.py} with alternative seeds.

\subsection{External Validity}

The evaluation covers two benchmark domains: an autonomous vehicle CPS
(9~states, 13~transitions) and two e-commerce models (24 and 42
states).  Both are well-documented, representative use cases for MBST,
but neither is a large-scale industrial system.  Scalability to models
with hundreds of states, or to domains with highly ambiguous
requirements (e.g., informal agile user stories), has not been
empirically validated.  The hierarchical Markov chain component
(Section~\ref{subsec:hierarchical}) addresses state-space scalability
architecturally but was not independently evaluated in this study.

The LLM oracle uses Azure OpenAI GPT-4.1-mini.  Substituting a
different model or a smaller open-weight model may change the oracle
resolution rate and downstream F1 scores.

\subsection{Statistical Validity}

All quantitative results derive from a deterministic proof-of-concept
implementation rather than a statistical sample of runs.  No
confidence intervals are reported.  Coverage results (100\% in
Table~\ref{tab:operational_metrics}) are guaranteed by the test
generator's design and do not require statistical treatment; F1 and
divergence metrics are point estimates from single runs on fixed
benchmarks.  Future work should replicate results across multiple
requirement documents and LLM providers to establish statistical bounds.
